%% ============================================================================
%% §2 Model and Numerical Design  --  joint
%%   2.1 MobileFaceNet topology
%%   2.2 Precision selection (PyTorch float32 vs C float32/INT16/INT8)
%%   2.3 Model evaluation (face-verification cosine similarity)
%% ============================================================================
\section{Model and Numerical Design}
\label{sec:model}

Before any hardware is built, the model topology and the numerical format must be fixed. Both decisions are constraints that the hardware datapath is then specialised around.

\subsection{MobileFaceNet Topology}
\label{sec:model:topology}

We target the MobileFaceNet variant with input size $3\times112\times96$. The original MobileFaceNet paper~\cite{chen2018mobilefacenetsefficientcnnsaccurate} and the Eyeriss-style accelerator in~\cite{7738524} use a square $3\times112\times112$ input with a $7\times7$ global depthwise convolution. We instead adopt the non-square variant from~\cite{Linjun_SUN20202019EDL8103}, which replaces the GDConv kernel with $7\times6$ to match the reduced horizontal resolution. The complete layer schedule is shown in Table~\ref{tab:mobile}. When each Bottleneck is expanded into its three constituent sub-layers (PW1, DW, PW2), the full model comprises \textbf{50 layers} (L0--L49).

\begin{table}[H]
\centering
\caption{MobileFaceNet Layer Configuration ($3\times112\times96$ input)}
\label{tab:mobile}
\begin{tabular}{lllrrrr}
\toprule
\textbf{Input shape} & \textbf{Operator} & & \textbf{t} & \textbf{c} & \textbf{n} & \textbf{s} \\
\midrule
$3\times112\times96$   & Conv $3\times3$        & & --  & 64  & 1 & 2 \\
$64\times56\times48$   & DW Conv $3\times3$     & & --  & 64  & 1 & 1 \\
$64\times56\times48$   & Bottleneck             & & 2   & 64  & 5 & 2 \\
$64\times28\times24$   & Bottleneck             & & 4   & 128 & 1 & 2 \\
$128\times14\times12$  & Bottleneck             & & 2   & 128 & 6 & 1 \\
$128\times14\times12$  & Bottleneck             & & 4   & 128 & 1 & 2 \\
$128\times7\times6$    & Bottleneck             & & 2   & 128 & 2 & 1 \\
$128\times7\times6$    & Conv $1\times1$        & & --  & 512 & 1 & 1 \\
$512\times7\times6$    & GDConv $7\times6$      & & --  & 512 & 1 & 1 \\
$512\times1\times1$    & Linear (PW $1\times1$) & & --  & 128 & 1 & 1 \\
\bottomrule
\end{tabular}
\end{table}

\subsection{Precision Selection}
\label{sec:model:precision}

A face-verification network must preserve its embedding's discriminative property under quantisation: the gap between intra-person and inter-person cosine similarity must remain large. We evaluated four precision formats (PyTorch float32, a plain-C float32 reference, C INT16 Q10, and C INT8 Q4) by computing pairwise cosine similarities on a face-verification task and comparing the largest absolute error against the PyTorch reference. Results are in Table~\ref{tab:precision}.

\begin{table}[H]
\centering
\caption{Precision Comparison Summary}
\label{tab:precision}
\begin{tabular}{lcc}
\toprule
\textbf{Format} & \textbf{Max error vs PyTorch} & \textbf{Verdict} \\
\midrule
C float32 & $\approx 0$                                      & Reference \\
C INT16   & $<0.003$ (most pairs: OK)                        & \textbf{Selected} \\
C INT8    & up to $0.34$ (many pairs: WARN)                  & Rejected \\
\bottomrule
\end{tabular}
\end{table}

The C float32 path closely matches the PyTorch reference, confirming the floating-point implementation. C INT16 (Q10 fixed-point) introduces only negligible error on the majority of pairs, while C INT8 has significantly larger deviations, especially on pairs whose similarity is already near zero (the regime that decides match/no-match). We therefore select \textbf{INT16 Q10} as the hardware representation: best trade-off between hardware cost and numerical accuracy for this workload.

Table~\ref{tab:fixed_point} fixes the bit-widths used downstream by all hardware modules.

\begin{table}[H]
\centering
\caption{Fixed-Point Format Selected for Hardware}
\label{tab:fixed_point}
\begin{tabular}{lll}
\toprule
\textbf{Quantity} & \textbf{Format} & \textbf{Width} \\
\midrule
Activation     & Q10 & \texttt{int16} \\
Weight         & Q10 & \texttt{int16} \\
Bias           & Q20 & \texttt{int32} \\
Accumulator    & Q20 & \texttt{int40} \\
PReLU $\alpha$ & Q10 & \texttt{int16} \\
\bottomrule
\end{tabular}
\end{table}

\subsection{Model Evaluation}
\label{sec:model:eval}

To validate the selected INT16 Q10 model end-to-end, intra-person and inter-person cosine similarities were measured on two datasets: a 6-person / 3-images-each public set and a self-prepared set containing four well-known individuals (Dan, Elon, Trump, Wil). Results in Table~\ref{tab:face_eval} show a large separability gap (intra $\gg$ inter) on both datasets, confirming that INT16 Q10 quantisation preserves the model's discriminative ability with negligible loss against the floating-point reference.

\begin{table}[H]
\centering
\caption{Face-Verification Evaluation (C INT16 Q10)}
\label{tab:face_eval}
\begin{tabular}{lccc}
\toprule
\textbf{Dataset} & \textbf{Intra avg} & \textbf{Inter avg} & \textbf{Gap} \\
\midrule
Original (6 persons, 3 images each)       & 0.5575 & 0.0673 & $+0.4903$ \\
Self-prepared (Dan, Elon, Trump, Wil)     & 0.5270 & 0.0783 & $+0.4487$ \\
\bottomrule
\end{tabular}
\end{table}

Finally, Table~\ref{tab:sim_compare} compares the C fixed-point reference against the RTL output on the self-prepared dataset. The two are bit-identical, confirming that the hardware datapath faithfully implements the chosen fixed-point arithmetic.

\begin{table}[H]
\centering
\caption{Similarity Summary: C Fixed-Point vs RTL}
\label{tab:sim_compare}
\begin{tabular}{lcc}
\toprule
\textbf{Metric} & \textbf{C fixed} & \textbf{RTL} \\
\midrule
Intra avg & 0.5270 & 0.5270 \\
Inter avg & 0.0783 & 0.0783 \\
Gap       & $+0.4487$ & $+0.4487$ \\
\bottomrule
\end{tabular}
\end{table}

This bit-exactness against the C-model golden is the foundation for all subsequent verification at the RTL, synthesised, and post-route levels (Section~\ref{sec:verify}).
